\section{Інструкція користувача}

\subsection{Встановлення та конфігурація}

\begin{lstlisting}[language=bash, caption={Встановлення залежностей}]
cd bridge-swarm-iot
make install          # python3 -m venv .venv && pip install -r requirements.txt
make up-full          # docker compose --profile observability up -d
\end{lstlisting}

Ключові параметри у файлі \texttt{.env}:
\begin{lstlisting}[caption={Фрагмент .env}]
SIM_NUM_DRONES=3
SIM_CRUISE_SPEED_MPS=0.5
SIM_TAKEOFF_ALTITUDE_M=3.0
SIM_BATTERY_DRAIN_RATE_PPS=0.05      # реалістичний: ~33 хв до 0%
SIM_BATTERY_RELAY_THRESHOLD_PCT=40.0
MQTT_HOST=localhost
MQTT_PORT=1883
\end{lstlisting}

\subsection{Запуск без Isaac Sim (headless)}

\begin{lstlisting}[language=bash, caption={Headless-демо естафетної передачі}]
# Термінал 1 — симулятор (прискорений drain для демо)
SIM_BATTERY_DRAIN_RATE_PPS=2.0 python -m swarm.mock_swarm

# Термінал 2 — запуск місії
python missions/relay_demo.py

# Відкрити дашборд
make grafana    # --> http://localhost:3000
\end{lstlisting}

\subsection{Запуск в Isaac Sim}

\begin{lstlisting}[language=bash, caption={Запуск з Pegasus Simulator}]
# у директорії Isaac Sim
./python.sh /шлях/до/bridge-swarm-iot/isaac/swarm_app.py

# в іншому терміналі
python missions/relay_demo.py
\end{lstlisting}

\subsection{Команди оператора (CLI)}

\begin{lstlisting}[language=bash, caption={Команди оператора}]
# Зліт / навігація
python -m swarm.operator takeoff --all --altitude 3.0
python -m swarm.operator goto uav-01 --x -3.25 --y 0.0 --z 3.0

# Інспекція
python -m swarm.operator inspect-poi uav-01 rack_left_end
python -m swarm.operator orbit uav-01 --cx -3.25 --cy 0 \
    --altitude 3.0 --radius 1.5

# Сканування коридора
python -m swarm.operator scan-zone aisle_left   # uav-01

# Зупинка / посадка
python -m swarm.operator stop --all
python -m swarm.operator land --all
\end{lstlisting}

\subsection{Критерії верифікації}

\begin{table}[h]
  \centering
  \caption{Критерії перевірки роботи системи}
  \begin{tabularx}{\linewidth}{lXl}
    \toprule
    Критерій & Перевірка & Очікуваний результат \\
    \midrule
    Зліт & \texttt{takeoff --all --altitude 3.0} & Усі 3 дрони на $z=3{,}0\pm0{,}1$~м \\
    Прямолінійний прохід & \texttt{scan-zone aisle\_left} & $|x - (-3{,}25)| < 0{,}1$~м вздовж всього Y \\
    Орбіта & \texttt{inspect-poi uav-01 rack\_left\_end} & Радіус $1{,}5\pm0{,}1$~м навколо POI \\
    Батарея & Grafana «Battery \%» panel & Лінійне зниження від 100\% \\
    Естафета & \texttt{relay\_demo.py}, drain=2\%/с & Handoff $< 0{,}5$~м від зупинки A \\
    Покриття & Grafana «Coverage \%» panel & Досягає 100\% після завершення проходу \\
    \bottomrule
  \end{tabularx}
  \label{tab:criteria}
\end{table}

\subsection{Одиничні тести}

\begin{lstlisting}[language=bash, caption={Запуск тестів}]
cd bridge-swarm-iot && python -m pytest tests/ -q
\end{lstlisting}

Результат: 49 тестів пройдено (0 помилок).

\begin{table}[h]
  \centering
  \caption{Одиничні тести по модулях}
  \begin{tabularx}{\linewidth}{llX}
    \toprule
    Файл тестів & К-сть & Що перевіряється \\
    \midrule
    \texttt{test\_commands.py}   & 9  & Round-trip всіх 7 команд; відхилення некоректних payload \\
    \texttt{test\_trajectory.py} & 12 & Квінтичний поліном: нульові $v, a$; \texttt{build\_orbit}: точки на колі, замкненість \\
    \texttt{test\_coverage.py}   & 8  & Boustrophedon: кількість waypoints, чергування напрямку \\
    \texttt{test\_facility.py}   & 10 & POI-реєстр, ScanZone non-overlap, altitude $\geq 3$~м \\
    \texttt{test\_missions.py}   & 6  & default\_pois, default\_zones; ScanCmd валідація \\
    \texttt{test\_operator.py}   & 4  & CLI публікує правильні MQTT-повідомлення \\
    \bottomrule
  \end{tabularx}
  \label{tab:tests}
\end{table}
